% ECONOMICS_PROBLEM: {"id": "E52-157", "title": "Monetary policy and inequality", "jel_code": "E52", "source_id": "OP-0249"}
% !TeX program = xelatex
\documentclass[11pt,a4paper]{article}
\usepackage[margin=23mm]{geometry}
\usepackage{fontspec}
\setmainfont{Latin Modern Roman}
\usepackage{amsmath,amssymb,mathtools}
\usepackage{xcolor,enumitem,booktabs,longtable,array,xurl,hyperref}
\definecolor{muted}{HTML}{53636C}
\hypersetup{colorlinks=true,urlcolor=blue,linkcolor=blue}
\setlength{\parindent}{0pt}
\setlength{\parskip}{5pt}
\newcommand{\R}{\mathbb R}
\newcommand{\E}{\mathbb E}
\newcommand{\N}{\mathbb N}
\newcommand{\ind}{\mathbf 1}
\newcommand{\Var}{\operatorname{Var}}
\newcommand{\Cov}{\operatorname{Cov}}
\newcommand{\supp}{\operatorname{supp}}
\DeclareMathOperator*{\argmin}{arg\,min}
\DeclareMathOperator*{\argmax}{arg\,max}
\newcommand{\entrymeta}[3]{{\sffamily\small\color{muted}\textbf{\texttt{#1}}\quad|\quad #2\par #3\par}}
\newcommand{\Problemref}[1]{\texttt{#1}}
\newcommand{\JELref}[2]{\texttt{#2} (JEL \texttt{#1})}
\begin{document}
\section*{Monetary policy and inequality}
\textbf{Problem ID:} \texttt{E52-157}\quad \textbf{JEL:} \texttt{E52}\par
% BEGIN ECONOMICS STATEMENT
\label{problem:OP-0249}
{\sffamily\small\color{muted}Original subject group: Inflation and monetary policy\par}
\label{definitions:problem:30007530}
\textbf{Audit classification:} Explicit published open question. Optimal HANK policy is expressly challenging; the finite-horizon, fixed-fiscal-rule problem is editorial. Verification concerns the cited version, not first priority or proof that this exact editorial specification remains unsolved on October 8, 2026.
\subsection*{Definitions and precise research target}
\paragraph{Editorial mathematical specification.} Fix a finite-horizon heterogeneous-agent monetary economy with incomplete insurance, idiosyncratic earnings risk, nominal assets, sticky prices, and horizon \(T\). Household types \(i\) have utility \(u_i(c,\ell)\), social weights \(\omega_i\ge0\), and discount factor \(\beta\). Initial distributions and a terminal resource and debt rule are given. Government transfers and taxes follow a fixed feasible rule, so that the monetary authority controls only an adapted policy-rate rule \(a=(a_0,\ldots,a_T)\) subject to an effective lower bound and an information constraint.
For every feasible adapted rule inducing an equilibrium, define
\[
W(a)=E\sum_{t=0}^T\beta^t\sum_i\omega_i u_i(c_{it}(a),\ell_{it}(a)).
\]
Characterize a maximizing rule and the marginal welfare terms associated with nominal redistribution, unemployment risk, price dispersion, and constrained consumption. Determine whether an implementable rule based on inflation, employment, and measured distributional states approximates this rule within a declared welfare tolerance. Identify the exposure and preference parameters needed for those comparisons, and report sensitivity to social weights rather than treating them as estimated facts. The finite horizon and fixed fiscal rule avoid assuming the existence of an unrestricted infinite-horizon Ramsey steady state. The research target is a quantitatively disciplined equity–efficiency policy assessment, not a claim that reducing any inequality statistic always improves monetary policy.

A stochastic policy is a sequence of decision functions \(a_t:\mathcal H_t\to[\underline i,\overline i]\), where \(\mathcal H_t\) is the specified observation history; a deterministic rate path is only the open-loop special case. Define a feasible set by all household budgets, firm optimality, nominal pricing constraints, market clearing, the fixed fiscal rule, and terminal debt repayment. Define policy approximation by \(W(a^*)-W(\widehat a)\le\epsilon_W\), using the same initial law and terminal condition. Inequality may be reported through the consumption distribution or a specified Gini index, but its reduction is not substituted for the declared household-utility objective. Welfare derivatives require differentiability; with occasionally binding constraints use finite policy contrasts instead.
\subsection*{Variants and assumption changes}
The following are \emph{editorial variants}, not additional published open conjectures.
\begin{enumerate}
\item \emph{Fixed fiscal redistribution.} Permit only the monetary rate rule to change and hold household transfers and terminal taxes fixed as functions of states. Quantify nominal redistribution and employment effects under these institutional limits.
\item \emph{Joint monetary-fiscal control.} Add a budget-feasible transfer instrument observed by the same authority. Compare the attainable welfare frontier and monetary rule with the fixed-fiscal benchmark; this is a different policy problem, not another solution to the restricted one.
\end{enumerate}
\subsection*{References and source audit}
\begin{itemize}
\item Adrien Auclert; Matthew Rognlie; Ludwig Straub. \emph{Fiscal and Monetary Policy with Heterogeneous Agents}. 2025. Locator: Section5, Optimal policy, printed pp.21--22; difficult quantitative welfare policy computation. Primary text checked. \url{https://web.stanford.edu/~aauclert/annual_review_hank.pdf}.
\end{itemize}
Optimal HANK policy is expressly challenging; the finite-horizon, fixed-fiscal-rule problem is editorial.
\begin{enumerate}\raggedright
\item\hypertarget{definitions:source:30007530:1}{} Adrien Auclert; Matthew Rognlie; Ludwig Straub. 2025. \emph{Fiscal and Monetary Policy with Heterogeneous Agents}. Section5, Optimal policy, printed pp.21--22; difficult quantitative welfare policy computation.
\url{https://web.stanford.edu/~aauclert/annual_review_hank.pdf}.
DOI: \url{https://doi.org/10.1146/annurev-economics-091624-044646}.
\end{enumerate}

\subsection*{Common conventions: Source mathematical conventions}

These conventions apply to each entry unless explicitly replaced.

\textbf{Models and identification.} A primitive model \(m\) specifies all probability laws, feasible choices, information, equilibrium and measurement rules used by its target. The admissible class \(\mathcal M\) is fixed independently of outcomes. If \(P_O(m)\) is its observable probability law and \(\tau(m)\) the target, its identified set at observable law \(P\) is
\[
 \mathcal I_\tau(P)=\{\tau(m):m\in\mathcal M,\ P_O(m)=P\}.
\]
Point identification means that this set is a singleton. An empty set signals incompatibility of data and assumptions. Coordinate bounds need not describe the joint set. Bounds are sharp only if every retained target is attainable or arbitrarily approachable by compatible models. Finite bounds require explicit restrictions on unobserved quantities; unrestricted confounding, hidden wealth, extrapolation or arbitrary equilibrium selection can yield uninformative sets.

\textbf{Causal interventions.} A potential outcome \(Y_i(p)\) is the outcome under a complete policy regime \(p\), including timing, financing, information and market spillovers. The target population and units are fixed before treatment. With interference, use the complete assignment vector or a justified exposure mapping; individual no-interference notation alone cannot identify an economy-wide intervention. A difference of mean potential outcomes does not require the cross-world joint distribution. Conditioning on post-treatment migration, survival, enrollment or employment changes the estimand and needs separate justification.

\textbf{Statistical statements.} An estimator, confidence set, forecast or test specifies a sampling law, model class and loss/coverage criterion. Uniform \((1-\alpha)\)-coverage means \(\inf_{m\in\mathcal M}\Pr_m\{\tau(m)\in C_n\}\geq1-\alpha\), or an explicitly asymptotic version. The whole parameter space trivially has coverage, so informativeness requires a stated width, risk or power objective. A forecasting improvement is relative to a named benchmark, information set and evaluation population.

\textbf{Equilibrium and optimization.} An equilibrium comprises feasible plans, optimality, beliefs where needed, budget constraints and all market/resource clearing. The primitives or a stated benchmark define each correspondence \(\mathcal E(p;m)\). Nonemptiness is assumed or proved, never inferred just from compact policy sets. A selected-equilibrium target uses a declared selection rule; a robust target quantifies over all permitted equilibria. Use suprema or infima unless attainment is established. An \(\varepsilon\)-optimal policy is feasible and within \(\varepsilon\) of the finite optimum value. Report an empty feasible set separately.

\textbf{Welfare and accounting.} Normative utility, interpersonal normalization, social weights, numeraire, discounting and the use of public receipts are primitives. Behavioral utility need not coincide with the welfare criterion. Household welfare includes ownership income and rents once; adding profits again to compensating variation generally double-counts them. A monetary equivalent is defined by an explicit transfer and price convention, with monotonicity and range conditions. Average effects, mechanism contributions, transfers and welfare losses are different targets. Infinite sums require convergence and dynamic feasible plans require terminal or no-Ponzi conditions.

\textbf{Source versions.} A working-paper date, revision date and journal publication year are distinguished. A DOI for a journal version is not automatically the DOI of an earlier manuscript. Locators refer to the stated version and printed pages where specified. A metadata-only check does not verify a claim about a full-text conclusion. Every entry's source subsection narrows the provenance claim accordingly.
% END ECONOMICS STATEMENT
\end{document}
